\chapter{Error Reporting}
\label{sec:errors}

A rule that rejects correct programs lives or dies by its messages. beni inherits Elm's standard for
them, and this chapter designs the ownership error to that standard.

\section{What a message must do}
\label{sec:errors-principles}

A message names \emph{three places}, states \emph{a consequence}, writes \emph{a fix}, and gives
\emph{a verdict} on whether the rejection is real sharing or a limit of the checker. The three places
are Rust's: the borrow, the invalidating action, and the later use, with the blame on the action---``the
error, conceptually, is to perform an invalidating action in between two uses of the reference''
\citep{rfc2094}. Futhark names the two sites, and its error index offers a copy as the fix on every one of its consumption errors, and notes
when a copy ``subverts the purpose of using in-place updates in the first place''
\citep{futhark2026errorindex}; its authors also learned that unnamed intermediate results must be
described in terms the user can read, as in ``Consuming result of applying "iota" (at 2:23-28)''
\citep{henriksen2021anatomy}. Hylo's compiler ``will offer to insert these copies for you''
\citep{hylo2024bindings}. Czaplicki's rules for Elm apply unchanged: the error shows the code exactly as
the programmer wrote it, and every message has a useful hint \citep{czaplicki2015errors}. Wrenn and
Krishnamurthi's view of an error message as a classifier of code into ``look here'' and ``don't look
there'' \citep[\S2]{wrenn2017classifiers}, with a term relevant when an edit there could resolve the
error, fixes the highlight set: the kept place, the edit, the later use, nothing else. Marceau, Fisler
and Krishnamurthi warn that highlighting one side of an inconsistency has ``an over-focusing effect''
and that ``expected \ldots\ found'' wording ``suggests that the definition is fine and the use
problematic'' \citep[\S7.2]{marceau2011mind}; that is why the message blames none of the three places
alone.

Crichton, Gray and Krishnamurthi's study of Rust learners shapes the verdict line. Participants ``could
usually predict why the borrow checker would reject a program'' (64--78\,\% of the time), but ``could only fix the program in
46\,\% of cases \ldots\ and could only create a counterexample in 31\,\% of cases''; asked about one
rejected program that was in fact safe, only 3 of 15 participants identified that it was
\citep[\S2.4]{crichton2023grounded}. As Crichton put it earlier, a sound and incomplete analysis's
errors ``may arise from either ownership-unsound behavior or limitations of the analyzer.
Understanding this distinction is essential for fixing ownership errors''
\citep{crichton2020usability}. Every message therefore says which it is, and how the compiler knows.

\section{Examples}
\label{sec:errors-examples}

\paragraph{A later read (real sharing).}
\begin{Verbatim}
-- EDIT OF A LIST STILL IN USE ----------------------------------- src/Main.beni

`List.push` changes a list in place, but this one is still needed afterwards.

12|     saved = model.todos
                ^^^^^^^^^^^ `saved` keeps the list here
15|     next = List.push model.todos todo
               ^^^^^^^^^^^^^^^^^^^^^^^^^^ the list is changed in place here
19|     List.length saved
                    ^^^^^ and `saved` is read here, after the change

If the change happened in place, `saved` would already contain `todo` when it
is read on line 19: `List.length saved` would be one more than before.

To keep `saved` as it was, copy it:

    saved = List.copy model.todos

Or, if `saved` only needs the list before the change, use it before line 15.

This is sharing the checker can see: the list is read on line 19 on every path
through this function.
\end{Verbatim}

\paragraph{A closure that keeps the list (real sharing).}
\begin{Verbatim}
`List.set` changes a list in place, but a function still holds it.

 8|     show = λi → String.fromInt (List.length rows)
                                                ^^^^ `show` holds `rows` here
11|     rows2 = List.set rows 0 r
                ^^^^^^^^^^^^^^^^^ the list is changed in place here
14|     { model | label = show 3 }
                          ^^^^ `show` is called here, after the change

`show` would see the changed list. If that is intended, call `show` with a
copy: `show = λi → String.fromInt (List.length (List.copy rows))` — or compute
what `show` needs before line 11.
\end{Verbatim}

\paragraph{A command that captured the list (real sharing, across messages).}
\begin{Verbatim}
`List.push` changes `model.todos` in place, but a command started by another
message may still be reading it.

-- in the arm for `Save`:
41|     Cmd.perform (λsend → send (Saved (encode model.todos)))
                                                 ^^^^^^^^^^^ the command keeps `model.todos`
-- in the arm for `Add`:
27|     { model | todos = List.push model.todos todo }
                          ^^^^^^^^^^^^^^^^^^^^^^^^^^ changed in place here

A command runs after `update` returns and may run after later messages, so the
list it keeps is the same one `Add` changes. To give the command its own copy:

41|     Cmd.perform (λsend → send (Saved (encode (List.copy model.todos))))

Or encode before the command is made:

40|     text = encode model.todos
41|     Cmd.perform (λsend → send (Saved text))
\end{Verbatim}

\paragraph{A closure that consumes a capture, called many times (real sharing).}
\begin{Verbatim}
This function changes `acc` in place, but `List.map` may call it once per element.

 7|     List.map xs (λx → List.push acc x)
                     ^^^^^^^^^^^^^^^^^^^^ changes `acc` in place
        ^^^^^^^^ `List.map` calls it for every element of `xs`

Each call would change the same list, so the result of the first call would
change under the second. To build a list by adding to `acc`, fold:

 7|     List.foldl xs acc (λx acc → List.push acc x)
\end{Verbatim}

\paragraph{A limit of the checker.}
\begin{Verbatim}
`List.set` changes `rows` in place, but the checker cannot tell whether anything
else still holds it.

14|     rows = Js.to raw
               ^^^^^^^^^ `rows` came from JavaScript here, which may keep it
17|     List.set rows i v
        ^^^^^^^^^^^^^^^^^ changed in place here

This is a limit of the checker, not proof of sharing: nothing in this module
reads the old `rows` after line 17, but the checker cannot see what JavaScript
does with the array it gave you. To be safe, copy it once where it arrives:

14|     rows = List.copy (Js.to raw)
\end{Verbatim}

\section{The fixed shape}
\label{sec:errors-shape}

Every message has the same skeleton:
\begin{Verbatim}
<TITLE> ------------------------------------------------------------ <file>

<one sentence: what primitive, what list, why that is a problem here>

<the kept place, in source order, with a label beginning "… keeps"/"… holds">
<the edit, with the label "changed in place here">
<the later use, with the label "… read/called here, after the change">

<one or two sentences: what the program would see if the change happened in place>

<the fix, machine-applicable, as a replacement line, "To … , copy it: …">
<an alternative fix that avoids the copy, when there is one>

<the verdict line: "This is sharing the checker can see: …" or "This is a limit
 of the checker, not proof of sharing: …", naming the limit>
\end{Verbatim}
That is: a title; one sentence naming the primitive, the list and why that is a problem here; the kept
place, the edit and the later use, in source order, labelled ``keeps'' or ``holds'', ``changed in place
here'', and ``read'' or ``called here, after the change''; one or two sentences on what the program would
see if the change happened in place; the fix as a replacement line, with an alternative that avoids the
copy where one exists; and the verdict line. The suggested copy is mechanically applicable in rustc's
sense \citep{rustc2026diagnostics}: the compiler can insert \code{List.copy} at the kept place by itself.
The message does not say ``did you mean'', and it does not call the program ``illegal''.

The design ships the fix and the verdict line, against one recorded dissent: Marceau, Fisler and
Krishnamurthi's advice, from their study of novices, that error messages should not propose
solutions \citep{marceau2011mind}. Hylo, Futhark, Rust and Swift all ship fixes, and an ownership
error's fix is mechanical in a way that a type error's is not, which is why we follow them.

\section{Telling a limit from real sharing}
\label{sec:errors-tags}

Every fact the checker derives carries a \emph{derivation tag}: \emph{exact} when it comes from
something the source shows---a variable occurrence, a field access, a call whose summary was inferred
from a beni body, a capture in a lambda written here---or \emph{approximate} when it comes from one of
the checker's own limits, listed in Table~\ref{tab:tags}.

\begin{table}[t]
\centering\small
\begin{tabularx}{\linewidth}{@{}>{\raggedright\arraybackslash}p{0.42\linewidth}X@{}}
\toprule
Limit & Text of the tag \\
\midrule
a \code{foreign} or \code{Js} operation declared shared, or with no summary & ``came from / was given to JavaScript, which may keep it'' \\
a cell made $\top$ by the path cut or a width cap & ``the value is nested deeper than the checker follows (8 steps)'' \\
a path-insensitive join & ``the list may be either \code{a} or \code{b} here; the checker treats it as both'' \\
a read from a \code{Ref} or \code{Queue} & ``was taken from a \code{Ref}, whose other readers the checker cannot see'' \\
a pessimistic summary on an unannotated platform \code{foreign} & ``\code{M.f} declares nothing about what it keeps'' \\
$\mathit{calls} = \mathit{many}$ for a closure the program calls once in fact & ``\code{List.map} may call it more than once'' \\
a container whose contents are not exclusive & ``the elements of \code{rows} may share'' \\
\bottomrule
\end{tabularx}
\caption{Approximate derivation tags and their wording.}
\label{tab:tags}
\end{table}

A rejection whose derivation uses only exact tags is \emph{real sharing}: there is a path through the
program text on which the old version is read after the edit. One that uses an approximate tag is
\emph{possible sharing}, and the message names the limit. This makes Crichton's distinction mechanical,
and it is also what the evaluation counts (\S\ref{sec:evaluation}). The tags give a third thing for free:
a rejection with no derivation at all is an internal compiler error, never a user error---the analogue
of Swift's move-checker diagnostics that report ``a compiler bug'' when a use cannot be verified
\citep{swift2026diagnosticssil}.

\section{Stability}
\label{sec:errors-stability}

\begin{quote}
A program the checker of release $n$ accepts is accepted by every release $m > n$.
\end{quote}

The accepted set is defined by the rules of \S\ref{sec:model}--\S\ref{sec:inference}, never by what an
optimiser removes. In a 2025 discussion of an explicit-copies mode for Swift, McCall argued that such
guarantees ``would have to be real guarantees, and the compiler would be responsible for doing them even
in debug builds'', and Farvardin that the checking ``needs to happen earlier in the optimization
pipeline, just to guarantee stability across different versions of the compiler''
\citep{swift2025explicitcopies}. Precision improvements are additive---a finer path language, a
container function that gains an exact summary, a closure position that gains $\mathit{calls} \le
1$---and each only removes approximate tags. The one legitimate exception is a \emph{soundness fix}: a
rule found unsound may reject previously accepted programs, and must do so with a diagnostic that says
so and, where a mechanical rewrite exists, a migration in the formatter. The rule set is versioned in
the interface format, so a cached interface produced under an older rule set is a cache miss, never a
silent mix.
